\styledchapter{Interdependent-Values Correlated Private Information}

Suppose there is one unit of an indivisible good for sale with $N$ bidders and $1$ seller. Bidders do not know their values. Instead, each bidder receives a signal $s_i$ about her value $v_i$. A bidder's signal is a random variable before it is received, but then deterministic once it is realized.

Let uppercase letters denote random variables and lowercase letter denote their realization. We have the random variables $V_1,\ldots,V_N,S_1,\ldots,S_N$ with joint density $	f(v_1,\ldots,v_N, s_1,\ldots,s_N)$.

\section{Symmetry and Affiliation}

\begin{assumption}[Symmetry]
	We assume that for every $i \ne j$ and $\left( v,s \right)= \left( v_1,\ldots,v_N, s_1,\ldots,s_N \right)$, \[
		f(v_1,\ldots,v_N, s_1,\ldots,s_N) = f(v_1',\ldots,v_N',s_1',\ldots,s_N')
	\] where $(v',s') = (v_1',\ldots,v_N',s_1',\ldots,s_N')$ is obtained by substituting $v_i, s_i$ for $v_j,s_j$ from $(v,s)$.

	Along this line, we assume that all bidders use the same bidding function.
\end{assumption}

\begin{assumption}[Signals and values]
	We require $s_i,v_i \in [0,\infty)$.
\end{assumption}

Immediately, this implies \[
	\E\left[V_1 \mid S_1 = s_1, S_2=  s_2, \ldots, S_N = s_N\right]  = \E\left[V_1 \mid S_1 = s_1, S_2 = s_2', \ldots,S_N = s_N'\right] 
,\] where $(s')$ is a permutation of $(s)$.\boxednote{Note that $s_1$ does not get swapped with anything, because in the definition of symmetry, if we swap $s_i$ with $s_j$, we must also swap $v_i$ with $v_j$.}

\begin{remark}	The correlated private information auction is a generalization of the independent private values model; we can impose $s_i = v_i$.
\end{remark}

\begin{assumption}[Affiliation]
	Assume that for all $i$, $\overline{v} = \left( \overline{v}_1,\ldots,\overline{v}_N \right) \ge \left( \underline{v_1},\ldots,\underline{v_N} \right) = \overline{s}$, and $\overline{s} = \left( \overline{s}_1,\ldots,\overline{s}_N \right) \ge \left( \underline{s}_1,\ldots,\underline{s}_N \right) = \underline{s}$,\footnote{Here, the ordering on vectors is defined by the same ordering on all components.} we have $
	\frac{f\left( v_i, \overline{v}_{-i}, \overline{s}_{-i} \right)}{f\left( v_i, \underline{v}_{-i}, \underline{s}_{-i} \right)}$ is nondecreasing in $v_i$ on $[0,\infty)$ and $\frac{f(\overline{v}, s_i, \overline{s}_{-i})}{f(\underline{v}, s_{i}, \underline{s}_{-i})}$ is nondecreasing in $s_i$ on $[0,\infty)$ and strictly increasing if $\overline{v}_i > \underline{v}_i$.
\end{assumption}

\begin{proposition}	Given these assumptions of symmetry and affiliation,
		\begin{enumerate}[label=(\roman*)]
			\item $\E\left[V_1 \mid S_1=s_1,S_2= s_2, \ldots,S_N = s_N\right] $ is nondecreasing in $s_2,\ldots,s_N$ and strictly increasing in $s_1$.
			\item $\E\left[\phi(S_1,\ldots,S_N) \mid a_1 \le S_1 \le b_1,\ldots,a_N \le S_N \le b_N\right]$ is nondecreasing in $a_1,b_1,\ldots,a_N,b_N$ whenever $\phi$ is nondecreasing in each of its arguments.
		\end{enumerate}
	\end{proposition}

\begin{assumption}[Equilibrium bidding function]
	Given symmetry, we also assume that all bidders use the same equilibrium bidding strategy.\boxednote{This rules out asymmetric equilibria, which certainly exist.} Furthermore, we assume that bidding is more aggressive when signals are higher. So, one's most serious competitor is the bidder among the others with the highest signal.
\end{assumption}

\begin{notation}	We can focus on bidder $1$ without loss.

	Let $Y_1,Y_2,\ldots,Y_{N-1}$ denote the highest, second-highest, and so on of the bidders' signals $S_2,\ldots,S_N$.
	Let $X_1,X_2,\ldots,X_{N-1}$ denote the highest, second-highest, and so on of all the signals $S_1,.S_2,\ldots,S_N$.
\end{notation}

\section{Second-Price Auction}
Suppose bidder 1 receives a signal $s_1$. How is she to bid?
We conjecture that bidder 1 should bid $\E\left[V_1 \mid S_1 = s_1\right] $.

Actually, considering the winner's curse\footnote{If the bid $\E\left[V_i \mid S_i = s_i\right] $ is nondecreasing in $s_i$, the winner is likely to have gotten a high signal that overestimates the value of the good.}, bidding $\E\left[V_1 \mid S_1 = s_1\right] $ is naive and non-strategic.

\begin{definition}[Winner's curse]
	Suppose in a second price auction that all bidders bid according to $\tilde{b}(s) = \E\left[V_1 \mid S_1 = s\right]$, which is strictly increasing by affiliation in $s$. Then the bidder with the highest signal wins. If your signal is $s$, then $\E\left[V_1 \mid S_1 = s\right] = \E\left[V_1 \mid S_1 = s, 0 \le S_j < \infty \forall\  j \ne 1\right]$. Conditioning on $s$ being the signal that wins, \begin{align*}
		\E\left[V_1 \mid S_1 = s = X_1\right] &= \E\left[V_1 \mid S_1 = s,\ 0 \le S_j < s \forall\  j \ne 1\right] \\
													 &\le \E\left[V_1 \mid S_1 = s_1, \ 0 \le S_j < \infty \forall\  j\ne 1\right] \\
													 &= \E\left[V_1 \mid S_1 = s\right] 
	,\end{align*} 
	where the inequality is strict if there is actually correlation (usually the case).

	There is typically a gap between your bid $\tilde{b}(s)$ and what the good is worth to you conditional on winning.
\end{definition}

If the second highest bid is between this gap, the winner of a SPA would pay more than her value if she bid according to $\E\left[V_1 \mid S_1 = s_1\right] $

Let $b^*(s)$ denote the bidding function that all bidders use and suppose $b^*(s)$ is strictly increasing. Assuming that $b^*$ is an equilibrium function, let's find $b^*$ (and later, check that it actually constitutes and equilibrium).

	Note that $b^*(y_1)$ is the price bidder $1$ pays if she wins.
	Consider
	\[
		\E\left[V_1 \mid S_1 = s, Y_1 = y_1\right],
	\]
	the expected value conditional on winning.
	Suppose this has a higher intercept and lower slope than $b^*(y_1)$ due to affiliation.

Then you should bid $b^*(s)$, which is the value of the intersection point of the two functions.
We see that $b^*(s) = \E\left[V_1 \mid S_1 = s, Y_1 = s\right] $.

\begin{remark}	Strategic bidding is different from estimating the value of the good given your signal: \[
		\E\left[V_1 \mid S_1 = s, Y_1 = s\right] \ne \E\left[V_1 \mid S_1 = s\right] 
	.\] 
	Suppose $s>y_1$, which means that $b^*(s)$ is a wining bid. Then 
	\begin{align*}
		b^*(s) = \E\left[V_1 \mid S_1 = s, Y_1 = s\right]
		&\ge \underbrace{\E\left[V_1 \mid S_1= s, Y_1 = y_1\right]}_{\text{correct value estimate}}  \\ &> \E\left[V_1 \mid S_1 = y_1, Y_1 =y_1\right] \\ &= b^*(y_1) .\end{align*} There is no winner's curse because the correct value estimate is strictly greater than the price paid. (If a bidder wins, she does not regret the bid ex-post.) Now suppose $s<y_1$, so $b^*(s)$ is a losing bid. The \begin{align*} b^*(s) = \E\left[V_1 \mid S_1 = s, Y_1 = s\right] &\le \underbrace{\E\left[V_1 \mid S_1 =s, Y_1 = y_1\right]}_{\text{correct value estimate}} \\ &< \E\left[V_1 \mid S_1 = y_1, Y_1 =y_1\right] \\
		&= b^*(y_1)
	.\end{align*}
	There is no loser's curse because the correct value estimate is less than the price required to win.

	Winning bids overestimate the true value, while losing bids underestimate the true value.
\end{remark}

Suppose the seller also has private information, namely, a signal $S_0$. Should the seller reveal or conceal her information? How much should she reveal?

\begin{proposition}	If $X$ and $Z$ have joint density $f(x,z)$, then \[
		\E_X\left[x \mid Z =z\right] = \int_{}^{} x f(x \mid z)\,dx 
	.\] This implies \[
		\E_Z\left[\E_X\left[X \mid Z\right] \right] = \E\left[X\right] 
	.\] 
\end{proposition}

Assume $S_0,S_1,S_2,\ldots,S_n,V_1,\ldots,V_N$ are affiliated and $S_1,S_2,\ldots,S_N$ are ex ante symmetric. For now, consider two possibilities: either the seller commits to always revealing $S_0$ or to always conceal $S_0$. If the seller commits to never revealing her information, then the equilibrium is as before: $b^*(s) = \E\left[V_1 \mid S_1 = s, Y_1= s\right] $.

If she commits to always revealing her information, then
\[
	b^*(s,s_0) = \E\left[V_1 \mid S_1 = s, Y_1 = s, S_0 = s_0\right].
\]

By the Law of Iterated Expectations, \[
	\E\left[V_1 \mid S_1 = s, Y_1 = s\right] = \E_{S_0}\left[\E\left[V_1 \mid S_1 = s, Y_1=s, S_0\right] \mid S_1 = s, Y_1 = s\right] 
,\] so \[
	b^*(s) = \E_{S_0}\left[b^*(s,S_0) \mid S_1 = s, Y_1 = s\right] 
.\] 

Let $R_U$ be the seller's expected revenue when bidders are always uninformed and $R_I$ when they are always informed.
Let $X_1 \ge X_2 \ge \cdots \ge X_N$ be the order statistics of $S_1,.S_2,\ldots,S_N$. Then \begin{align*}
	R_U &= \E\left[b^*(X_2)\right] \\
	R_I &= \E\left[b^*\left( X_2, S_0 \right)\right] \\
		&= \E_{X_2}\left[\E_{S_0}\left[b^*(S_2,S_0) \mid X_2\right]  \right]\\
	R_I - R_U &= \E_{X_2}\left[\E_{S_0}\left[b^*(X_2,S_0) \mid X_2\right] - b^*(X_2)\right] 
.\end{align*} 
We conjecture this value is nonnegative.
To show it, it suffices to prove that fixing any value $X_2 = s$ guarantees \[
	b^*(s) \le \E_{S_0}\left[b^*\left( s,S_0 \right)\mid X_2 = s\right] 
.\] We see that
\begin{align*}
	b^*(s) &=  \E_{S_0}\left[b^*(s,S_0) \mid S_1=s, Y_1=s\right]  \\
		   &= \E_{S_0}\left[b^*(s,S_0) \mid S_1 = X_2 = s, Y_1 = s\right] \text{ since $Y$ are the values of the \emph{others}} \\
		   &\le \E_{S_0}\left[b^*(s,S_0) \mid S_1=X_2=s, Y_1 \ge s\right] \\
		   &= \E_{S_0}\left[b^*(s,S_0) \mid S_1=X_2 =s\right] \text{ because $\P\left[Y_1 = s\right] = 0$} \\
			   &= \E_{S_0}\left[b^*\left( s,S_0 \right)\mid X_2=s\right]
	.\end{align*}
	where the last step uses symmetry: conditioning on $S_1=s$ provides no additional information once $X_2=s$ is fixed.

For intuition on why the seller should reveal/conceal, first consider the concealment case. If the seller's revenue comes from bidder $1$, who bids \[
	b^*(s) = \E\left[V_1 \mid S_1 = s, Y_1 = s\right] 
,\] then this bid being a losing bid (it's the second-highest bid) implies that the $Y_1 > s$ and thus $b^*(s)$ actually underestimates the true value $V_1$.
A more accurate estimate of $V_1$ would be $\E\left[V_1 \mid S_1 = s, Y_1 \ge s\right]  > b^*(s)$.
By revealing $S_0,$, the seller induces bidder $1$ to bid $b^*(s,S_0) = \E\left[V_1\mid S_1= s, Y_1 = s, S_0\right] .$
Because $S_0$ is affiliated with $Y_1$ and $V_1$, conditioning on $S_0$ will on average help to correct the bias in the estimate of $V_1$. For losing bids, the bias is downwards, so revealing $S_0$ tends to increase losing bids (and in particular, the second-highest bid).\boxednote{We can conjecture that in a FPA, the seller should \emph{not} reveal her information.}

Suppose the seller can now \emph{sometimes} reveal her information. We only consider that the case where the bidder reveals if the signal is above $\alpha$ and conceals if it is below. The bidders also know this information, as otherwise, they may think that the seller is always revealing/concealing, and this implies the seller is deceiving the bidders. We won't consider this as a possibility!

By mirroring our previous analysis, when the seller gets a low signal and conceals, the bidders know the signal is low and recondition. But then the seller should reveal her information to raise the second-highest bid, so assuming the seller and bidders have the same understanding of the mechanism, the seller's optimal policy is to always reveal her information.

\section{English Auction}
We will suppose by symmery that all bidders use the same bidding strategy that can be described by the $N-1$ functions $\left( b^*_0, b_1^*,\ldots,b_{N-2}^{*} \right)$, where given your signal $s$, $b_0^*(s)$ is the price you drop out if no one has yet dropped out, $b_1^*(s \mid p_1)$ is the price at which you drop out when $1$ bidder has dropped out at $p_1$, and so on, and $b_k^*(s \mid p_1,\ldots,p_k)$ is the price at which you drop out when $k$ bidders have dropped out at $p_1 < p_2 < \cdots < p_k$.

Let's derive these $b_k^*$ functions starting from $b_0^*$.

Like before, first suppose that $\left( b_0^*,\ldots,b_{N-2}^{*} \right)$ is a symmetric equilibrium and each $b_k^*(s \mid p_1,\ldots,p_k)$ is strictly increasing in $s$ with the intention to verify it later.

Suppose that all others drop out simultaneously. In this case, they must have the same signal. This is the only scenario that makes $b_0^*$ relevant. Let $y := y_1=y_2= \cdots = y_N$. 
Then $b^*_0(s) = \E\left[V_1 \mid S_1 = s, Y_1 = Y_2 = \cdots = Y_N = y\right]$. Note that $b_0^*(s)$ is strictly increasing.

Suppose one bidder has dropped out first at price $p_1$. We defined $b_1^*(s \mid p_1)$ as the price at which you drop out when $1$ bidder has dropped out at price $p_1$. 
Observing $p_1$ allows all remaining bidders to infer that $Y_{N-1} = y_{N-1}$, where $p_0^*(y_{N-1}) = p_1$.

If all remaining bidders have the same signal $y$, then \[
	b_1^*(s \mid p_1) = \E\left[V_1 \mid S_1 = s, Y_1 = \cdots = Y_{N-2} = s, Y_{N-1} = y_{N-1}\right] 
.\] While $p_1$ does not explicitly appear on the right hand side, $y_{N-1}$ is determined by $p_1$ through $b_0^*(y_{N-1}) = p_1$. We can continue this pattern of proof by graph and find that in general,
\[
	\begin{aligned}
		b_k^*(s\mid p_1,\ldots,p_k)
		&= \E\left[V_1 \mid S_1=  s, Y_1 = \cdots = Y_{N-k-1} = s,\right. \\
		&\qquad \left.Y_{N-k} = y_{N-k}, \ldots,Y_{N-1} = y_{N-1}\right].
	\end{aligned}
\]
Here $b_0^*(y_{N-1}) = p_1$, $b_1^*(y_{N-2}) = p_2$, $\ldots$, and $b_{k-1}^*(y_{N-k} \mid p_1,\ldots,p_{k-1}) = p_k$.

Does your drop out price ever jump below the clock price when someone else drops out?
\begin{proposition}	A bidder's drop out price never jumps below the clock price when someone else drops out.
\end{proposition}
\begin{proof}
	We begin by providing intuition for a particular case. Suppose $y_{N-1} < s$.
	\begin{align*}
		p_1 = b_0^*(y_{N-1}) &= \E\left[V_1 \mid S_1 = Y_1 - \cdots = Y_{n-1} = y_{N-1}\right] \\
							 &< \E\left[V_1 \mid S_1 = Y_1 = \cdots = Y_{N-2}, Y_{N-1} = y_{N-1} \right] = b_1^*(s_1 \mid p_1) \\
							 &\le \E\left[V_1 \mid S_1=Y_1 =\cdots = Y_{N-1} = s\right] \\
							 &= b_0^*(s)
	.\end{align*}

	Therefore, $p_1 < b_1^*(s \mid p_1) \le b_0^*(s)$.

		In general, if $y_{N-1} < y_{N-2} < \cdots < y_{N-k} < s$, then
		\begin{align*}
			p_k = b_{k-1}^*\left( y_{N-k} \mid p_1,\ldots,p_{k-1} \right)
			&= \E\left[V_1 \mid S_1 = Y_1 = \cdots = Y_{N-k} = y_{N-k},\right. \\
			&\qquad \left.\underbrace{Y_{N_k+1} = y_{N-k+1}, \ldots, Y_{N-1} = y_{N-1}}_{\text{$k-1$ known signals}}\right]  \\
			&< \E\left[V_1 \mid S_1 = Y_1 \cdots  = Y_{N-k-1} = s,\right. \\
			&\qquad \left.Y_{N-k} = y_{N-k}, \ldots,Y_{N-1} = y_{N-1}\right]  \\
			&= b_k^*(s \mid p_1,\ldots,p_k) \\
			&\le \E\left[V_1 \mid S_1 = Y_1 =\cdots= Y_{N-k} = s,\right. \\
			&\qquad \left.Y_{N-k+1}= y_{N-k+1}, \ldots,Y_{N-1} = y_{N-1}\right] \\
			&= b_k^{*}(s \mid p_1,\ldots,p_k) \\
			&= b_{k-1}^*(s \mid p_1,\ldots,p_{k-1})
		.\end{align*} 
\end{proof}

\begin{proposition}	$\left( b_0^*,b_1^*,\ldots,b_{N-2}^* \right)$ is a symmetric equilibrium.
\end{proposition}

\begin{proof}
	We must show that if all others use the $b_k^*$ functions to set their drop-out prices, then you can do no better than do so as well.

	Suppose your signal is $s$. Let the others' signals be $y_1 \ge y_2 \ge \cdots \ge y_{N-1}$. Suppose the second-highest signal is $y_{N-1} < s$.
	To win, you have to pay \begin{align*}
		b_{N-2}^*\left( y_1 \mid p_1,p_2,\ldots,p_{N-2} \right) 
		&= \E\left[V_1 \mid S_1 = Y_1 = y_1, Y_2 = y_2, \ldots,Y_{N-1} = y_{N-1}\right]
	.\end{align*} 

		What is the good worth to you given $y_1 \ge y_2 \ge \cdots \ge y_{N-1}$?
		It is worth the first term in
		\[
			\begin{aligned}
				\E\left[V_1 \mid S_1 = s_1, Y_1 = y_1,\ldots,Y_{N-1} = y_{N-1}\right]
				&\gtreqless
				\E\left[V_1 \mid S_1 = y_1,\ldots,Y_{N-1} = y_{N-1} \right] \\
				&\iff s \gtreqless y_1
			\end{aligned}
		.\]
		Hence, you want to win exactly when $s > y_1$.
		You can do this by using the bidding strategy, since drop out prices are increasing in the signal.
\end{proof}

\section{Revenue Comparison between the Second-Price Auction and English Auction}

We conjecture the ascending price will raise more revenue than the second-price auction, since it reveals more information. (Recall that the seller would ex ante always prefer to reveal information she has; so it would make sense for the seller to like others to reveal information as well.)

Let $X_1 \ge X_2 \ge \cdots \ge X_N$ be the order statistics of the $N$ signals $S_1,\ldots,S_N$.

Let $\SPA_2$ be a second price actuion in which the seller observes all but the two highest signals and reveals the signals she observes.
So, she observes and reveals $X_3=x_2,\ldots,X_N = x_N$, then allows the bidders with the highest signals to bid in a canoncial SPA. The equilibrium bidding strategy is \[
	b_{\SPA_2}^{*}\left( s \mid x_3,\ldots,x_N \right)= \E\left[V_1 \mid S_1=Y_1=s, X_3= x_3,\ldots,X_n = x_n\right] 
.\] 

Therefore, if $X_1=  x_1 \ge X_2 = x_2 \ge \cdots \ge X_N = x_N$, then \[
	R_{\SPA_2}(x_1,\ldots,x_N) = \E\left[V_1 \mid S_1 = Y_1 = x_2, Y_2 = x_3, \ldots,Y_{N-1}=x_N\right] 
.\] In this situation, if we had run an ascending auction instead, then the revenue would be the same. This is because the final drop-out price of the person with the second-highest signal is determined by conditioning on all lower signals. Thus, \[
	R_E\left( x_1,\ldots,x_N \right) = R_{\SPA_2}(x_1,\ldots,x_N) \ge R_{\SPA}(x_1,\ldots,x_N)	
,\]

\begin{conj}
	Second-price auctions are used in the real world because of risk aversion.
\end{conj}

While the highest signal wins in these auctions, the auctions are not necessarily efficient, because someone with a high value can draw a low signal.

\section{First-Price Auction}
Let $f_{Y_1}(y\mid s)$ be the probability densiity that the highest signal of others is $y$ given that your signal is $s$. By affiliation, \[
	\frac{f_{Y_1}(\overline{y} \mid s)}{f_{Y_1}(y \mid s)}
\] is nondecreasing in $s$ for every $\overline{y} > y$. 

This implies that
\begin{equation} \label{2-12-0}
	\frac{f_{Y_1}(y \mid s)}{F_{Y_1}(y \mid s)}
\end{equation}
is nondecreasing in $s$ for every $y$.

Define $w(s,y) = \E\left[V_1 \mid S_1=s, Y_1 = y\right] $ and let $\hat{b}(\cdot )$ be a strictly increasing symmetric equilibrium bidding function (if it exists). Let $u(r,s)$ be your expected utility when you signal is $s$, you bid $\hat{b}(r)$, and all others bid according to $\hat{b}(\cdot )$. Then, as long as $\hat{b}(\cdot )$ is indeed strictly increasing (which, as before, we plan to check later), 

When the highest signal of the others is $y$, your expected utility is \begin{equation} \label{2-12-1}
	u(r,s) = \int_{0}^{r} \left( w(s,y) - \hat{b}(r) \right)f_{Y_1}(y\mid s) dy 
\end{equation} because $\hat{b}$ being strictly increasing implies that $\hat{b}(y) < \hat{b}(r)$ iff $y < r$. 
Thus, since $\hat{b}$ is an equilibrium, the FOC in $r$ of $u(r,s)$ implies \begin{align*}
	0 = \left. \frac{d}{dr}u(r,s) \right|_{r = s} &= \left. \left( w(s,r) - \hat{b}(r) \right)f_{Y_1} (r \mid s) + \int_{0}^{r} \left( -\hat{b}'(r) \right)f_{Y_1}(y\mid s) dy \right|_{r=s}  \\
											  &= \left.\left( w(s,r) - \hat{b}(r) \right) f_{Y_1}(r \mid s) - \hat{b}'(r) F_{Y_1}(r\mid s) \right|_{r=s} \\
												  \implies \hat{b}'(s) F_{Y_1}(s\mid s) &= \left( w(s,s) - \hat{b}(s) \right)f_{Y_1}(s\mid s)
	.\end{align*} 
The LHS is the marginal cost of increasing your bid, while the RHS is the marginal benefit. We can rearrange to se that for every signal $s$,
\begin{equation} \label{2-12-2}
	\hat{b}'(s) = \left( w(s,s) - \hat{b}(s) \right) \frac{f_{Y_1} (s\mid s)}{F_{Y_1}(s\mid s)}
\end{equation}
We can solve this ODE subject to the boundary condition $\hat{b}(0) = w(0,0)$:\footnote{We impose this boundary condition because it makes sense over the more sensible $\hat{b}\left( 0 \right) = 0$. For instance, we can give a circular reason. \texthl{fix.}}
\[
  \hat{b}'(s) + \frac{f_Y(s\mid s)}{F_{Y_1}(s\mid s)}\,\hat{b}(s)
  = w(s,s)\frac{f_{Y_1}(s\mid s)}{F_{Y_1}(s\mid s)},
  \qquad \hat{b}(0)=w(0,0).
\]

Define the integrating factor
\[
  \mu(s):=\exp\left(\int_{0}^{s}\frac{f_Y(x\mid x)}{F_{Y_1}(x\mid x)}\,dx\right).
\]
Multiplying the ODE by \(\mu(s)\) gives
\[
  \frac{d}{ds}\left(\mu(s)\hat{b}(s)\right)
  = \mu(s)\,w(s,s)\frac{f_{Y_1}(s\mid s)}{F_{Y_1}(s\mid s)}.
\]
Integrating from \(0\) to \(s\) and using \(\mu(0)=1\) yields
\[
  \mu(s)\hat{b}(s) - \hat{b}(0)
  = \int_{0}^{s}\mu(y)\,w(y,y)\frac{f_{Y_1}(y\mid y)}{F_{Y_1}(y\mid y)}\,dy,
\]
so \[\hat{b}(s)
  = w(0,0)\exp\left(-\int_{0}^{s}\frac{f_Y(x\mid x)}{F_{Y_1}(x\mid x)}\,dx\right)
+ \int_{0}^{s} w(y,y)\,h(y\mid s)dy.\]
Evaluating the integral $\int_\epsilon^{s} \frac{f_Y(x\mid x)}{F_{Y_1}(x\mid x)} dx$, we see that as $\epsilon \downarrow 0$, the integral diverges to $\infty$, so the first term in $\hat{b}$ is $0$.
Therefore,
\begin{equation} \label{2-12-3}
	\hat{b}(s) = \int_{0}^{s} w(y,y) h(y\mid s)dy 
\end{equation}
where
\[
  h(y\mid s)
  := \frac{f_{Y_1}(y\mid y)}{F_{Y_1}(y\mid y)}
     \exp\left(-\int_{y}^{s}\frac{f_Y(x\mid x)}{F_{Y_1}(x\mid x)}\,dx\right),
  \qquad 0\le y\le s.
\]
Letting $H(y \mid s) = \int_{0}^{y} h(z\mid s)dz $, we get \[
	H(y \mid s) = e^{-\int_{y}^{s} \frac{f_{Y_1}(x\mid x)}{F_{Y_1}(x\mid x)} dx }
\] for every $y \le s$. Thus, \[
H(s\mid s) = 1, \quad H(0\mid s) = 0
,\] and thus $h(y\mid s) > 0$ is a density on $y \in [0,s]$.

Equation (\ref{2-12-3}) is defined only for $s > 0$.
This equation says that $\hat{b}(s)$ is a weighted average of $w(y,y)$s with $y \in [0,s]$.
Since $w(y,y)$ is strictly increasing in $y$,
\begin{equation}
	\begin{aligned}
		\hat{b}(s)
		&= \int_{0}^{s} w(y,y) h(y \mid s) dy \\
		&< \int_{0}^{s} w(s,s) h(y\mid s) dy \\
		&= w(s,s) \int_{0}^{s} h(y\mid s )dy \\
		&= w(s,s)
	\end{aligned}
	\label{2-12-4}
\end{equation} 
for every $s$. 
Combining equations (\ref{2-12-2}) and (\ref{2-12-4}) yield that $\hat{b}'(s) > 0$ for every $s$. Therefore, (\ref{2-12-1}) is indeed the correct formula for $u(r,s)$ given $\hat{b}(\cdot )$ defined by (\ref{2-12-3}).

We have shown that $\hat{b}$ is increasing and that it satsfies the FOC of a maximization problem (by construction). It remains to check that it satisfies the second-order condition:
\begin{align*}
	\frac{d}{dr} u(r,s) &= \left( w(s,r) - \hat{b}(r) \right)f_{Y_1}(r \mid s) - \hat{b}'(R) F_{Y_1}(r \mid s) \\
						&= F_{Y_1}(r \mid s) \left[ w(s,r) - \hat{b}(r) \frac{f_{Y_1}(r \mid s)}{f_{Y_1}(r \mid s)} - \hat{b}'(r) \right] \gtreqless 0
.\end{align*}
Trying to analyze this in terms of fixing $r$ at $s$ and hypothesizing changes in $r$, fix $r=s$ and try changing $s$. By (\ref{2-12-2}) and $w(s,r)$ being strictly increasing in $s$ (both from affiliation), we see that decreasing $s$ to below $r$, the expression is negative, and increasing $s$ to above $r$, the expression is positive. Therefore, the derivative is maximized in $r$ at $r^* = s$.

\begin{aside}	What is the interpretation of $h(y\mid s)$? Reny and Brooks do not know, but Chen may.
\end{aside}

\begin{proposition}	We claim that $\frac{f_{Y_1}(y\mid s)}{F_{Y_1}(s\mid s)}$ is the density of second highest signal $X_2$ given that the highest signal $X_1$ is $s$.
\end{proposition}
\begin{proof}
	We know that $f_{Y_1}(y\mid s)$ is the density of $Y_1= y$ given $S_1 = s$. Thus, $\frac{f_{Y_1}(y\mid s)}{F_{Y_1}(a \mid s)}$ is the density of $Y_1 = y$ given $S_1=s$ and $Y_1 \le a$.
	We plug in $a=s$ to conclude, as \[
		Y_1 = y_1 \mid S_1 = s, Y_1 \le s \iff X_2 = y \mid S_1 = s = X_1 \iff X_2 = y \mid X_1 = s
	.\] 
\end{proof}

\begin{proposition}	The revenues of different auctions are ordered as \[
		R_\FPA \le R_\SPA \le R_E
	.\] 
\end{proposition}

\begin{proof}
	We have already shown the second inequality.

	In a second price auction, \[
		b^*(s) = \E\left[V_1 \mid S_1 = s, Y_1 = s\right] = w(s,s)
	.\] 
	Let $\left. R \right|_{X_1 = s} $ be the expected revenue given that the highest signal is $s$.
	We see that \[
		\left.R_\FPA\right|_{X_1 = s} = \hat{b}(s)
	,\] 
	while \[
		\left.R_\SPA\right|_{X_1 = s} = \int_0^{s}w(y,y) \P\left[X_2 = y \mid X_1 = s\right] dy= \int_{0}^{s} w(y,y) \frac{f_{y_1}(y\mid s)}{F_{Y_1}(s \mid s)} dy 
	.\] 
	If we can show that $\hat{b}(s)$ is less than this integral for every $s$, then we will have shown the revenue of an SPA is greater than that of a FPA.

	Recall that we found last time \[
		\hat{b}(s) = \int_{0}^{s} w(y,y) h(y\mid s)dy 
	.\] We want to show that $h(y \mid s)$ is greater than $\frac{f_{Y_1}(y\mid s)}{F_{Y_1}(s \mid s)}$ for large $y$.
	Hence, it suffices to show \[
		e^{- \int_{y}^{s} \frac{f_{Y_1}(x \mid x)}{F_{Y_1}(x \mid x)}dx } = H(y \mid s) \ge \frac{F_{Y_1} (y\mid s)}{F_{Y_1} (s \mid s)}
	\] for $y \le s$.
	We have that
	\begin{align*}
		H(y \mid s)
		&= \operatorname{exp}\left( - \int_{y}^{s} \frac{f_{Y_1} (x\mid x)}{F_{Y_1}(x\mid x)}dx \right) \\
		&\ge \operatorname{exp}\left(-\int_{y}^{s} \frac{F_{Y_1}(x\mid x)}{f_{Y_1}(x \mid x)}dx \right) \text{ since } \frac{f_{Y_1}(a \mid b)}{F_{Y_1}(a \mid b)} \text{ is nondecr. in $b$} \\
		&= \operatorname{exp}\left(- \left. \left( \ln{\left(F_{Y_1}(x \mid s)\right)}\right|_{y}^{s} \right)\right) \\
		&= \operatorname{exp}\left(- \left( \ln{\left(\frac{F_{Y_1} (s \mid s)}{F_{Y_1}(y \mid s)}\right)} \right)\right) \\
		&= \operatorname{exp}\left(\ln{\left(\frac{F_{Y_1} (y \mid s)}{F_{Y_1} (s \mid s)}\right)}\right) \\
		&= \frac{F_{Y_1} (y \mid s)}{F_{Y_1} (s \mid s)}
	.\end{align*}
\end{proof}
